\section{Note on sources and delivery}
This homily addresses an adult parish assembly for the resumed Fifth Sunday
after the Epiphany, General Roman Calendar, 8 November 2026, in the 1962
Roman Missal. It develops the reviewed study's interpretation
\texttt{patience-and-conversion}, with the compatible account of the
sinner's hope and Christ's resurrection from
\texttt{promise-prayer-and-pledge}. The relation of Gospel, Epistle and
orations is the homily's synthesis; neither Father is credited with an
interpretation of the complete Mass.

The owner's explicit concern for wheat injured by uprooting changes the
question from recognising wrongdoing to caring for persons endangered by
the proposed remedy. Augustine's distinction between plants and repentant
people explains why present wrongdoing does not exhaust a person's future.
Paul's paired commands to forgive and admonish specify active care during
the interval. The Secret includes the corrector among those needing pardon
and direction. The Postcommunion's pledge and sought effect provide the
conclusion in hope. Exact delivery figures and completed author proof
checks are recorded below and in \texttt{research/production-review.md}.

The spoken body contains 1,296 words. At an estimated unhurried
115--125 words per minute, delivery would take 10.4--11.3 minutes, before
additional pauses. This is a word-count estimate. The prose was reread
silently for sense and oral clarity; no audible rehearsal or timed human
delivery is claimed. The headings, this note, references and metadata
are outside the speech. No anecdote, personal experience, clerical
identity or newly composed recited prayer is supplied. The liturgical
source is the \work{Missale Romanum}, Vatican typical edition (1962).
The biblical exposition follows the \work{Douay--Rheims Bible},
Challoner revision, with the context distinguished from the appointed
extent in the canonical study.

\section{References}
\begin{itemize}
\item \work{Missale Romanum}, Vatican editio typica (1962),
pp.~420--421, nos.~1749--1758: Gospel, Epistle, Collect, Alleluia,
Offertory, Secret and Postcommunion; the canonical study retains the
collated Latin and identified historical English. Commentary here is
original prose, not a substitute translation of these prayers.
\item \work{Douay--Rheims Bible}, Challoner revision, composite
electronic witness, Gutenberg 1581: Matt 13:24--30,36--43;
Col 3:12--17; Ps 129:1--8 (Vulgate numbering; Hebrew 130).
The body uses the biblical scenes and arguments with few short phrases.
\item Augustine, Sermon XXIII [Benedictine LXXIII], §§1--4,
NPNF I.6; \work{Exposition of Psalm 129}, §§1--8,
NPNF I.8, heading Psalm 130. Retained historical English;
no fresh original-Latin collation is claimed.
\item John Chrysostom, \work{Homilies on Colossians}, VIII,
on Col 3:12--15, NPNF I.13, retained historical English:
forgiveness after serious injury, charity as the bond, and inward
affection beyond an outward concession.
\end{itemize}
The cited evidence and its qualifications belong to the reviewed canonical
research. Historical source expression retains its recorded public-domain
status and is excluded from the project's licence for original prose.
Author verification and independent homily review are separate events;
no external or ecclesiastical approval is claimed.
